\subsubsection{Análisis de Alternativas y Justificación - ¿Por qué pasamos de RUP a DAD? [Pendientes: 1. Completar comparación bibliográfica de ciclos y prácticas con las fuentes disponibles. 2. Definir y documentar las prácticas específicas de desarrollo que se utilizarán en el proyecto.]}\label{sec:spmp-8-1-1}
Rational Unified Process (RUP) fue considerado inicialmente por su naturaleza iterativa e incremental y su orientación a la arquitectura, características apropiadas para desarrollar la aplicación, construir la biblioteca de escenarios e integrar sus distintos componentes. Posteriormente se optó por DAD como una alternativa más reciente que conserva estas características pertinentes y permite adaptar la forma de trabajo a las necesidades del proyecto.

La elección de DAD responde a la necesidad de coordinar un producto de software con incertidumbres que no se resuelven únicamente mediante programación, puesto que la aplicación depende de decisiones sobre las entradas del simulador, su ejecución, el catálogo de escenarios y la aproximación espacial, por lo que revisar prioridades y obtener evidencia técnica desde las primeras iteraciones resulta coherente e importante para el trabajo previsto. La adaptación registrada conserva fases reconocibles y, al mismo tiempo, distribuye las pruebas, la arquitectura y la documentación entre las iteraciones; de este modo, la justificación de DAD frente a la alternativa inicialmente considerada se apoya en su adaptación al contexto y en la relación entre incertidumbre y resultados verificables~\cite{adaptacion2026}.
